% !TEX root = ../../../main.tex
% 所属：第二部分：多变量分析学 / 第3章：重积分
% 节：3.3 重积分的变量代换
% 正文按原文件顺序保存；由 ../chapter3.tex 调用。
\section{重积分的变量代换}
\subsection*{二重积分的变量代换}

\begin{definition}
	设 $U\subset\R^2$ 为开集，考虑映射
\[
T:U\longrightarrow V=T(U),\qquad
(u,v)\longmapsto\bigl(x(u,v),y(u,v)\bigr).
\]
要求 $T$ 为一一映射，并满足：
\begin{enumerate}
\item $x(u,v)$、$y(u,v)$ 在 $U$ 上连续可导；
\item Jacobian 行列式
\[
\frac{\partial(x,y)}{\partial(u,v)}
=\begin{vmatrix}
\dfrac{\partial x}{\partial u}&\dfrac{\partial x}{\partial v}\\[5pt]
\dfrac{\partial y}{\partial u}&\dfrac{\partial y}{\partial v}
\end{vmatrix}
\ne0.
\]
\end{enumerate}
这样的 $T$ 是 $U$ 到 $V$ 的一个微分同胚，称$T$为平面上的一个坐标变换。
\end{definition}

在 $U$ 中取直线 $u=u_0$，它在 $xy$ 平面中对应一条曲线
\[
x=x(u_0,v),\qquad y=y(u_0,v),
\]
称为 $v$-曲线。
同样，直线 $v=v_0$ 对应曲线
\[
x=x(u,v_0),\qquad y=y(u,v_0),
\]
称为 $u$-曲线。

由于 $T$ 为一一映射，$V$ 中的每一点既可唯一地用 $(x,y)$ 表示，
也可唯一地用 $(u,v)$ 表示。
称这些 $u$-曲线和 $v$-曲线构成曲线坐标网，
称 $(u,v)$ 为该点的曲线坐标。

\begin{center}
% 直线坐标网与曲线坐标网；固定 u 得到 v-曲线，固定 v 得到 u-曲线。
\begin{tikzpicture}[>=Stealth,font=\small]
  \begin{scope}
    \draw[->] (0,0) -- (3.15,0) node[right] {$u$};
    \draw[->] (0,0) -- (0,2.8) node[above] {$v$};
    \node[below left] at (0,0) {$O$};
    \foreach \u in {.85,1.45} \draw[gray] (\u,.25) -- (\u,2.4);
    \foreach \v in {.8,1.4} \draw[gray] (.3,\v) -- (2.65,\v);
    \node[above] at (.85,2.4) {$u=u_0$};
    \node[right] at (2.65,.8) {$v=v_0$};
  \end{scope}
  \draw[->] (3.55,1.45) -- (4.35,1.45) node[midway,above] {$T$};
  \begin{scope}[xshift=5cm]
    \draw[->] (0,0) -- (3.35,0) node[right] {$x$};
    \draw[->] (0,0) -- (0,2.8) node[above] {$y$};
    \node[below left] at (0,0) {$O$};
    \foreach \u in {.6,1.25,1.9}{
      \draw[domain=.3:2.15,samples=35,variable=\v]
        plot ({\u+.16*\v*\v},{\v+.08*\u*\u});
    }
    \foreach \v in {.65,1.25,1.85}{
      \draw[domain=.35:2.6,samples=35,variable=\u]
        plot ({\u+.16*\v*\v},{\v+.08*\u*\u});
    }
    \node[right] at (2.7,1.35) {$u$-曲线};
    \node[above] at (2.65,2.5) {$v$-曲线};
  \end{scope}
\end{tikzpicture}

\end{center}

\begin{theorem}[二重积分的变量代换]\label{p2-double-change-of-variables}
设 $U$ 是 $uv$ 平面上的开集，
\[
T:U\longrightarrow T(U),\qquad
(u,v)\longmapsto\bigl(x(u,v),y(u,v)\bigr)
\]
为一一映射。
设 $x(u,v)$、$y(u,v)$ 在 $U$ 上连续可导，且
$\partial(x,y)/\partial(u,v)\ne0$。
若 $D\subset U$ 为有界零边界闭区域，$f$ 在 $T(D)$ 上连续，则
\[
\iint_{T(D)}f(x,y)\,\dd x\,\dd y
=\iint_D f\bigl(x(u,v),y(u,v)\bigr)
\left|\frac{\partial(x,y)}{\partial(u,v)}\right|
\,\dd u\,\dd v.
\]
\end{theorem}
% 变量代换定理的补充证明见本节末。

\subsection*{变量代换的例子}

在下面的例子中，原来的 $xy$ 平面区域记为 $D$，
对应的参数平面区域记为 $D'$，即 $T(D')=D$。

\begin{example}[线性代换计算面积]
设 $a>0$，求曲线 $(x-y)^2+x^2=a^2$ 所围区域的面积。
\end{example}
\begin{solution}
原区域为
\[
D=\{(x,y):(x-y)^2+x^2\le a^2\},
\qquad
S_D=\iint_D\dd x\,\dd y.
\]
令
\[
u=x,\qquad v=x-y.
\]
则逆变换为
\[
x=u,\qquad y=u-v,
\]
并且
\[
\frac{\partial(x,y)}{\partial(u,v)}
=\begin{vmatrix}1&0\\1&-1\end{vmatrix}=-1.
\]
由于
\[
(x-y)^2+x^2=u^2+v^2,
\]
原区域对应圆盘
\[
D'=\{(u,v):u^2+v^2\le a^2\}.
\]
因此
\[
S_D=\iint_{D'}1\,\dd u\,\dd v=\pi a^2.
\]
\end{solution}

\begin{example}[双曲线与射线围成区域的面积]
设 $0<a<b$、$0<p<q$，求第一象限内由
\[
xy=p,\qquad xy=q,\qquad y=ax,\qquad y=bx
\]
围成的区域 $D$ 的面积。
\end{example}
\begin{solution}

\begin{center}
% xy=p、xy=q、y=ax、y=bx 围成的第一象限区域。
\begin{tikzpicture}[>=Stealth,font=\small,scale=1.05]
  \draw[->] (0,0) -- (3.75,0) node[right] {$x$};
  \draw[->] (0,0) -- (0,3.55) node[above] {$y$};
  \node[below left] at (0,0) {$O$};
  \fill[gray!12]
    (.816497,1.224745)
    plot[domain=.816497:1.490712,samples=35] (\x,{1/\x})
    -- (2.260777,1.017350)
    plot[domain=2.260777:1.238278,samples=35] (\x,{2.3/\x})
    -- cycle;
  \draw[domain=.35:3.4,samples=60] plot (\x,{1/\x});
  \draw[domain=.72:3.5,samples=60] plot (\x,{2.3/\x});
  \draw (0,0) -- (3.35,1.5075) node[right] {$y=ax$};
  \draw (0,0) -- (2.15,3.225) node[above] {$y=bx$};
  \node[right] at (3.42,.28) {$xy=p$};
  \node[right] at (3.5,.7) {$xy=q$};
  \node at (1.45,1.2) {$D$};
\end{tikzpicture}

\end{center}

令
\[
u=xy,\qquad v=\frac yx.
\]
由于 $x,y>0$，可解得
\[
x=\sqrt{\frac uv},\qquad y=\sqrt{uv}.
\]
原区域变为矩形
\[
D'=\{(u,v):p\le u\le q,\ a\le v\le b\}.
\]
计算 Jacobian 行列式：
\[
\begin{aligned}
\frac{\partial(x,y)}{\partial(u,v)}
&=\begin{vmatrix}
\dfrac1{2\sqrt{uv}}&-\dfrac12\sqrt{\dfrac u{v^3}}\\[7pt]
\dfrac12\sqrt{\dfrac vu}&\dfrac12\sqrt{\dfrac uv}
\end{vmatrix}\\
&=\frac1{4v}+\frac1{4v}
=\frac1{2v}.
\end{aligned}
\]
因此
\[
\begin{aligned}
S_D
&=\iint_D\dd x\,\dd y\\
&=\iint_{D'}\frac1{2v}\,\dd u\,\dd v\\
&=\frac12\int_p^q\dd u\int_a^b\frac1v\,\dd v\\
&=\frac12\int_p^q\ln\frac ba\,\dd u\\
&=\frac{q-p}{2}\ln\frac ba.
\end{aligned}
\]

\end{solution}

\subsection*{极坐标变换}

对于含有 $x^2+y^2$ 的表达式，可以考虑极坐标变换
\[
x=r\cos\theta,\qquad y=r\sin\theta.
\]
其 Jacobian 行列式为
\[
\frac{\partial(x,y)}{\partial(r,\theta)}
=\begin{vmatrix}
\cos\theta&-r\sin\theta\\
\sin\theta&r\cos\theta
\end{vmatrix}
=r.
\]
例如，圆域 $x^2+y^2<a^2$ 在参数平面中对应
$0<r<a$、$0<\theta<2\pi$，另需处理原点和角度接缝。
\textbf{半径约定：}本段用 $a$ 表示圆域半径，
所以写成 $x^2+y^2<a^2$，与径向上限 $r<a$ 相对应。

注意，当 $r=0$ 时，极坐标变换不是一一对应的。
在应用变量代换公式时，还需避开角度接缝，即在 $xy$ 平面中去掉非负实轴。
可以先在避开这些点的区域上应用公式，再取极限。

\begin{example}[圆域上的极坐标积分]
计算
\[
I=\iint_D\sin\!\left(\pi\sqrt{x^2+y^2}\right)\,\dd x\,\dd y,
\qquad
D=\{(x,y):x^2+y^2\le1\}.
\]
\end{example}
\begin{solution}
引入极坐标 $x=r\cos\theta$、$y=r\sin\theta$。
形式上，$D$ 对应参数区域
\[
D_1=\{(r,\theta):0\le r\le1,\ 0\le\theta\le2\pi\}.
\]
于是
\[
\begin{aligned}
&=\iint_{D_1}(\sin\pi r)\,r\,\dd r\,\dd\theta\\
&=\int_0^{2\pi}\dd\theta\int_0^1r\sin\pi r\,\dd r\\
&=2\pi\left[-\frac r\pi\cos\pi r
+\frac{\sin\pi r}{\pi^2}\right]_0^1\\
&=2.
\end{aligned}
\]

严格地说，需先避开原点和非负实轴。
对充分小的 $\varepsilon>0$，取参数矩形
\[
D_{1,\varepsilon}
=\{(r,\theta):\varepsilon\le r\le1,\
\varepsilon\le\theta\le2\pi-\varepsilon\},
\]
并记其在 $xy$ 平面中的像为 $D_\varepsilon$。
在这个参数矩形上，极坐标变换是一一的，且 Jacobian 行列式不为零。

\begin{center}
% 极坐标下避开原点与角度接缝：参数矩形对应去心扇环。
\begin{tikzpicture}[>=Stealth,font=\small]
  \begin{scope}
    \draw[->] (0,0) -- (2.8,0) node[right] {$r$};
    \draw[->] (0,0) -- (0,3.15) node[above] {$\theta$};
    \node[below left] at (0,0) {$O$};
    \draw[line width=.65pt] (.6,.55) rectangle (2.15,2.55);
    \draw[dashed] (0,.55) -- (.6,.55);
    \draw[dashed] (0,2.55) -- (.6,2.55);
    \draw[dashed] (.6,0) -- (.6,.55);
    \node[left] at (0,.55) {$\varepsilon$};
    \node[left] at (0,2.55) {$2\pi-\varepsilon$};
    \node[below] at (.6,0) {$\varepsilon$};
    \node[below] at (2.15,0) {$1$};
  \end{scope}
  \begin{scope}[xshift=5.7cm,yshift=1.4cm]
    \draw[->] (-1.65,0) -- (1.85,0) node[right] {$x$};
    \draw[->] (0,-1.65) -- (0,1.75) node[above] {$y$};
    \node[below left] at (0,0) {$O$};
    \draw[line width=.7pt]
      (18:1.3) arc[start angle=18,end angle=342,radius=1.3]
      -- (342:.23) arc[start angle=342,end angle=18,radius=.23]
      -- cycle;
  \end{scope}
\end{tikzpicture}

\end{center}

因此，按原图避开接缝后取极限的步骤，
\[
\begin{aligned}
I
&=\lim_{\varepsilon\to0^+}
\iint_{D_\varepsilon}
\sin\!\left(\pi\sqrt{x^2+y^2}\right)\,\dd x\,\dd y\\
&=\lim_{\varepsilon\to0^+}
\int_\varepsilon^{2\pi-\varepsilon}\dd\theta
\int_\varepsilon^1r\sin\pi r\,\dd r\\
&=\lim_{\varepsilon\to0^+}
(2\pi-2\varepsilon)
\left[-\frac r\pi\cos\pi r
+\frac{\sin\pi r}{\pi^2}\right]_\varepsilon^1\\
&=2.
\end{aligned}
\]

\textbf{补充说明：}记本例的被积函数为
$f(x,y)=\sin(\pi\sqrt{x^2+y^2})$。
上面取极限的依据是，$|f(x,y)|\le1$，
而删去的部分包含于半径为 $\varepsilon$ 的圆盘
以及总角度为 $2\varepsilon$ 的单位圆扇形之并。
所以
\[
\left|\iint_D f\,\dd x\,\dd y
-\iint_{D_\varepsilon}f\,\dd x\,\dd y\right|
\le m(D\setminus D_\varepsilon)
\le\pi\varepsilon^2+\varepsilon\longrightarrow0.
\]
\end{solution}

\subsubsection*{广义极坐标变换}

对于椭圆
\[
\frac{x^2}{a^2}+\frac{y^2}{b^2}\le1,
\qquad a>0,\ b>0,
\]
可取
\[
x=ar\cos\theta,\qquad y=br\sin\theta.
\]
相应的参数范围为 $0\le r\le1$、$0\le\theta\le2\pi$，
且
\[
\frac{\partial(x,y)}{\partial(r,\theta)}
=\begin{vmatrix}
a\cos\theta&-ar\sin\theta\\
b\sin\theta&br\cos\theta
\end{vmatrix}
=abr.
\]

\subsubsection*{平移后的极坐标变换}

例如，对于圆域
\[
(x-x_0)^2+(y-y_0)^2\le a^2,
\]
可先平移坐标，再作极坐标变换：
\[
x=x_0+r\cos\theta,\qquad
y=y_0+r\sin\theta.
\]
参数区域为 $0\le r\le a$、$0\le\theta\le2\pi$，
原点和角度接缝按前述方法处理。

\begin{center}
% 平移后的圆域与极坐标参数矩形。
\begin{tikzpicture}[>=Stealth,font=\small]
  \begin{scope}
    \draw[->] (0,0) -- (3.1,0) node[right] {$x$};
    \draw[->] (0,0) -- (0,2.8) node[above] {$y$};
    \node[below left] at (0,0) {$O$};
    \draw[line width=.65pt] (1.55,1.6) circle[radius=.7];
    \fill (1.55,1.6) circle[radius=.025];
    \node[below,font=\scriptsize] at (1.55,1.45) {$(x_0,y_0)$};
  \end{scope}
  \draw[->] (3.5,1.3) -- (4.3,1.3);
  \begin{scope}[xshift=5cm]
    \draw[->] (0,0) -- (2.9,0) node[right] {$r$};
    \draw[->] (0,0) -- (0,2.8) node[above] {$\theta$};
    \node[below left] at (0,0) {$O$};
    \draw[line width=.65pt] (0,0) rectangle (1.85,2.15);
    \node[left] at (0,2.15) {$2\pi$};
    \node[below] at (1.85,0) {$a$};
  \end{scope}
\end{tikzpicture}

\end{center}

\subsection*{变量代换定理的证明}



\subsubsection*{内外矩形逼近与零面积集的像}

取一个包含 $D$ 的闭矩形 $S$。
把它的两边不断二等分，得到一列逐次加细的矩形网格，网格的模长趋于零。
记完全包含于 $D$ 的小矩形之并为 $A_n$，
与 $D$ 相交的小矩形之并为 $B_n$，则
\[
A_n\subset A_{n+1}\subset D\subset B_{n+1}\subset B_n.
\]
由 $D$ 可求面积，
\[
mA_n\longrightarrow mD,\qquad mB_n\longrightarrow mD,
\qquad mB_n-mA_n\longrightarrow0.
\]
特别地，把与 $D$ 相交但不完全包含于 $D$ 的网格小矩形记为 $R_j$，则
\[
\sum_jmR_j=mB_n-mA_n\longrightarrow0.
\]
这些闭矩形只在零面积的边界上重叠，因此这里的面积相加不会重复计算正面积部分。
虽然 $S$ 不一定包含于 $U$，但 $D$ 是 $U$ 中的紧集，
所以充分细的网格中所有与 $D$ 相交的小矩形都包含于 $U$。

\subsubsection*{特殊变换下小矩形的面积}

\begin{lemma}[保留一个坐标的变换]\label{p2-triangular-area}
设 $T$ 是前述坐标变换，并且具有下列形式之一：
\[
T(u,v)=(u,y(u,v)),\qquad T(u,v)=(x(u,v),v).
\]
我们称这样的映射为\textbf{本原映射}.

若闭矩形 $R=[e,f]\times[h,j]$ 包含于 $U$，其中 $e<f$、$h<j$，
则存在 $(\widetilde u,\widetilde v)\in R$，使
\[
mT(R)=|J_T(\widetilde u,\widetilde v)|\,mR.
\]
\end{lemma}
\begin{proof}
先设 $T(u,v)=(u,y(u,v))$。此时 $J_T=y_v\ne0$。
由于 $R$ 连通，$y_v$ 在 $R$ 上符号不变。
先讨论 $y_v>0$ 的情形。

固定 $u$，函数 $v\mapsto y(u,v)$ 严格递增，故
\[
T(R)=\{(x,y):e\le x\le f,\ y(x,h)\le y\le y(x,j)\}.
\]

\[
mT(R)=\int_e^f\bigl(y(x,j)-y(x,h)\bigr)\,\dd x.
\]

\begin{center}
\begin{tikzpicture}[>=Stealth,font=\small,scale=.86]
  \draw[->] (0,0)--(2.8,0) node[right] {$u$};
  \draw[->] (0,0)--(0,2.4) node[above] {$v$};
  \draw (.7,.6) rectangle (2.1,1.8);
  \node[below] at (.7,0) {$e$};
  \node[below] at (2.1,0) {$f$};
  \node[left] at (0,.6) {$h$};
  \node[left] at (0,1.8) {$j$};
  \draw[dashed,gray] (.7,0)--(.7,.6) (2.1,0)--(2.1,.6)
    (0,.6)--(.7,.6) (0,1.8)--(.7,1.8);
  \node at (1.4,1.2) {$R$};
  \draw[->] (3,1.2)--(4.1,1.2) node[midway,above] {$T$};
  \begin{scope}[shift={(4.6,0)}]
    \draw[->] (0,0)--(3.6,0) node[right] {$x$};
    \draw[->] (0,0)--(0,2.7) node[above] {$y$};
    \draw (.7,.65) .. controls (1.1,1.2) and (1.6,.45) .. (2.1,.95);
    \draw (.7,1.8) .. controls (1.1,2.35) and (1.6,1.6) .. (2.1,2.1);
    \draw (.7,.65)--(.7,1.8) (2.1,.95)--(2.1,2.1);
    \draw[dashed,gray] (.7,0)--(.7,.65) (2.1,0)--(2.1,.95);
    \node[below] at (.7,0) {$e$};
    \node[below] at (2.1,0) {$f$};
    \node[right] at (2.1,.95) {$y(x,h)$};
    \node[right] at (2.1,2.1) {$y(x,j)$};
    \node at (1.4,1.45) {$T(R)$};
  \end{scope}
\end{tikzpicture}

\end{center}

先对上面的一元积分应用积分中值定理，再对 $v\mapsto y(\widetilde u,v)$
应用微分中值定理，得到 $\widetilde u\in[e,f]$、$\widetilde v\in(h,j)$，使
\begin{align*}
mT(R)
&=(f-e)\bigl(y(\widetilde u,j)-y(\widetilde u,h)\bigr)\\
&=(f-e)(j-h)y_v(\widetilde u,\widetilde v)\\
&=mR\,|J_T(\widetilde u,\widetilde v)|.
\end{align*}
若 $y_v<0$，交换上下两条边界曲线，便得到相同的绝对值公式。
对于 $T(u,v)=(x(u,v),v)$，交换两个坐标的作用，按同样的步骤证明即可。
\end{proof}

\subsubsection*{由小矩形的面积公式推出积分公式}

\begin{lemma}
	变量代换公式对本原映射成立.
\end{lemma}
\begin{proof}
	仍用第一步中的网格。记完全包含于 $D$ 的小矩形为
$D_1,\ldots,D_{N_n}$，其余与 $D$ 相交的小矩形为 $R_1,\ldots,R_{M_n}$。
令
\[
E_j=D\cap R_j.
\]
这些集合都是零边界紧集；它们及其像都可求面积。
因为 $T$ 是一一映射，它们的像也只在零面积边界上重叠。

对每个 $D_i$，按假设取 $(\widetilde u_i,\widetilde v_i)\in D_i$，使
\[
mT(D_i)=|J_T(\widetilde u_i,\widetilde v_i)|\,mD_i.
\]
再任取 $(a_j,b_j)\in E_j$。由 $T$ 在 $D$ 上一致连续，
当网格的模长趋于零时，上述各像集的直径也一致趋于零。
因此下面两个式子分别是 $f$ 在 $T(D)$ 上以及
$f(x(u,v),y(u,v))|J_T(u,v)|$ 在 $D$ 上的 Riemann 和：
\begin{align*}
S_n
&=\sum_{i=1}^{N_n}
 f\bigl(T(\widetilde u_i,\widetilde v_i)\bigr)mT(D_i)
 +\sum_{j=1}^{M_n}f\bigl(T(a_j,b_j)\bigr)mT(E_j),\\
S'_n
&=\sum_{i=1}^{N_n}
 f\bigl(T(\widetilde u_i,\widetilde v_i)\bigr)
 |J_T(\widetilde u_i,\widetilde v_i)|\,mD_i\\
&\quad+\sum_{j=1}^{M_n}
 f\bigl(T(a_j,b_j)\bigr)|J_T(a_j,b_j)|\,mE_j.
\end{align*}
这里 $T(u,v)$ 写在 $f$ 的自变量位置时，表示有序对
$\bigl(x(u,v),y(u,v)\bigr)$。

记这两个和式的第二部分分别为 $B_n$、$B'_n$。
由小矩形的面积公式，它们的第一部分完全相同，因而
\[
S_n-S'_n=B_n-B'_n.
\]
取常数 $H,K>0$，使 $|f|\le H$、$|J_T|\le K$。
有
\begin{align*}
|B_n|
&\le H\sum_jmT(E_j)
 \le HK\sum_jmR_j\longrightarrow0,\\
|B'_n|
&\le HK\sum_jmE_j
 \le HK\sum_jmR_j\longrightarrow0.
\end{align*}
两个被积函数都连续，而相应积分区域均为零边界紧集，所以它们都可积。
令网格模长趋于零，就得到
\begin{equation}\label{p2-triangular-substitution}
\iint_{T(D)}f(x,y)\,\dd x\,\dd y
=\iint_D f\bigl(x(u,v),y(u,v)\bigr)|J_T(u,v)|\,\dd u\,\dd v.
\end{equation}
\end{proof}



\subsubsection*{一般变换的局部分解}

\begin{lemma}[局部分解]\label{p2-local-triangular-factorization}
在定理 \ref{p2-double-change-of-variables} 的条件下，
对每个 $(u_0,v_0)\in U$，适当交换自变量或像坐标的次序后，
可以在 $(u_0,v_0)$ 的一个邻域内将 $T$ 分解为
\[
T=T_2\circ T_1,
\]
其中 $T_1$、$T_2$ 都是本原映射.
\end{lemma}
\begin{proof}
因为 $J_T(u_0,v_0)\ne0$，四个一阶偏导数中至少有一个非零。
适当交换坐标后，可设 $x_u(u_0,v_0)\ne0$。
由连续性，缩小邻域后可使 $x_u$ 在其中处处非零。

令
\[
T_1:\quad \xi=x(u,v),\qquad \eta=v.
\]
其 Jacobian 为 $x_u\ne0$。由反函数定理，进一步缩小邻域后，
$T_1$ 有连续可导的逆映射，可写为
\[
u=g(\xi,\eta),\qquad v=\eta.
\]
再令
\[
T_2:\quad x=\xi,\qquad y=y\bigl(g(\xi,\eta),\eta\bigr).
\]
这就得到 $T_2\circ T_1=T$。

对恒等式 $x(g(\xi,\eta),\eta)=\xi$ 求偏导，得到
\[
g_\xi=\frac1{x_u},\qquad g_\eta=-\frac{x_v}{x_u}.
\]
所以
\begin{align*}
\frac{\partial(x,y)}{\partial(\xi,\eta)}
&=\frac{\partial}{\partial\eta}
 y\bigl(g(\xi,\eta),\eta\bigr)\\
&=y_ug_\eta+y_v
 =\frac{x_uy_v-x_vy_u}{x_u}
 =\frac{J_T}{x_u}\ne0.
\end{align*}
上式右边的偏导数均在 $(u,v)=(g(\xi,\eta),\eta)$ 处取值。
$T_2=T\circ T_1^{-1}$ 也是一一映射，故两个映射都满足所需条件。
坐标交换的 Jacobian 的绝对值为 $1$，不会改变最后公式中的绝对值因子。
\end{proof}

设 $D'$ 是上述分解邻域内的任意零边界紧集。
先对 $T_2$ 使用已经证明的特殊情形，再对 $T_1$ 使用一次，得到
\begin{align*}
\iint_{T(D')}f(x,y)\,\dd x\,\dd y
&=\iint_{T_1(D')}
 f\bigl(T_2(\xi,\eta)\bigr)
 \left|\frac{\partial(x,y)}{\partial(\xi,\eta)}\right|
 \,\dd\xi\,\dd\eta\\
&=\iint_{D'} f\bigl(T(u,v)\bigr)
 \left|\frac{\partial(x,y)}{\partial(\xi,\eta)}
       \frac{\partial(\xi,\eta)}{\partial(u,v)}\right|
 \,\dd u\,\dd v\\
&=\iint_{D'}f\bigl(T(u,v)\bigr)|J_T(u,v)|\,\dd u\,\dd v.
\end{align*}
第二行中的第一个 Jacobian 要在 $(\xi,\eta)=T_1(u,v)$ 处取值；
最后一步使用了复合映射的 Jacobian 乘法公式。
这样就证明了变量代换公式的局部形式。

\subsubsection*{有限覆盖与整体公式}

\begin{proof}[定理 \ref{p2-double-change-of-variables} 的证明]
对每个 $p\in D$，按局部分解引理取开邻域 $U_p\subset U$，
使刚才的局部公式在 $U_p$ 内成立。
再取 $r_p>0$，使开球 $O(p,2r_p)\subset U_p$。
开球族 $\{O(p,r_p):p\in D\}$ 覆盖紧集 $D$。
由 Heine--Borel 定理，可取有限子覆盖
\[
D\subset\bigcup_{k=1}^sO(p_k,r_{p_k}).
\]
令
\[
\delta=\min_{1\le k\le s}r_{p_k}>0.
\]

用直径小于 $\delta$ 的网格小矩形 $Q_j$ 分割 $D$，记 $D_j=D\cap Q_j$，
忽略空集。
任取 $q\in D_j$，存在 $k$ 使 $q\in O(p_k,r_{p_k})$。
对任意 $q'\in D_j$，有
\[
|q'-p_k|\le|q'-q|+|q-p_k|<\delta+r_{p_k}\le2r_{p_k}.
\]
因此 $D_j\subset O(p_k,2r_{p_k})\subset U_{p_k}$，
可以在每个 $D_j$ 上使用局部公式。

各 $D_j$ 的公共部分都在网格分割线上，面积为零，
其像的面积也为零。
对有限个小块的局部公式求和，并利用积分的区域可加性，得到
\begin{align*}
\iint_{T(D)}f(x,y)\,\dd x\,\dd y
&=\sum_j\iint_{T(D_j)}f(x,y)\,\dd x\,\dd y\\
&=\sum_j\iint_{D_j}
 f\bigl(x(u,v),y(u,v)\bigr)|J_T(u,v)|\,\dd u\,\dd v\\
&=\iint_D f\bigl(x(u,v),y(u,v)\bigr)
 \left|\frac{\partial(x,y)}{\partial(u,v)}\right|\,\dd u\,\dd v.
\end{align*}
这就完成了定理的证明。
\end{proof}

\subsection*{高维情形与柱面、球面坐标}

\begin{theorem}[重积分的变量代换]
设 $U\subset\R^n$ 为开集，$T:U\to V=T(U)$ 为一一映射，
\[
T:\quad y_i=y_i(x_1,\ldots,x_n),\qquad i=1,\ldots,n.
\]
设各 $y_i$ 连续可导，且 Jacobian 行列式处处非零。
其逆映射也连续可导，记为
\[
x_i=x_i(y_1,\ldots,y_n),\qquad i=1,\ldots,n.
\]
由逆映射的链式法则，
\[
\frac{\partial(x_1,\ldots,x_n)}{\partial(y_1,\ldots,y_n)}\bigl(T(x)\bigr)
=\left(\frac{\partial(y_1,\ldots,y_n)}{\partial(x_1,\ldots,x_n)}(x)\right)^{-1}
\ne0.
\]
若 $\Omega\subset U$ 为有界零边界闭区域，$F$ 在 $T(\Omega)$ 上连续，则有
\[
\begin{aligned}
&\int\cdots\int_{T(\Omega)}F(y_1,\ldots,y_n)\,\dd y_1\cdots\dd y_n\\
&\quad=\int\cdots\int_\Omega
 F\bigl(y_1(x),\ldots,y_n(x)\bigr)
 \left|\frac{\partial(y_1,\ldots,y_n)}{\partial(x_1,\ldots,x_n)}\right|
 \,\dd x_1\cdots\dd x_n,
\end{aligned}
\]
其中 $x=(x_1,\ldots,x_n)$。
\end{theorem}
\begin{proof}
先考虑只改变一个坐标的变换
\[
H(t_1,\ldots,t_n)
=(t_1,\ldots,t_{k-1},h(t),t_{k+1},\ldots,t_n),
\qquad h_{t_k}\ne0.
\]
对闭长方体 $Q$，固定其余 $n-1$ 个坐标，
像集沿第 $k$ 个坐标的截面是由两个端点值确定的区间。
由累次积分公式，它的体积是这些区间长度在其余坐标上的积分。
先对这个连续的区间长度用积分中值定理，
再对第 $k$ 个坐标用一元微分中值定理，便得到
\[
mH(Q)=|h_{t_k}(\widetilde t)|\,mQ
=|\det DH(\widetilde t)|\,mQ
\]
对某个 $\widetilde t\in Q$ 成立。

将二维证明中的小正方形换成小正方体，
零体积集的像仍为零体积集：若 $H$ 在紧集附近满足 Lipschitz 估计，
边长为 $a$ 的小正方体的像可用体积不超过
$(2\sqrt nL)^na^n$ 的正方体覆盖。
因此原来的边界估计及 Riemann 和证明逐项适用，
从而得到这种单坐标变换的积分公式。

下面说明一般变换可以局部分解成这些变换。
固定一点，由 Jacobian 矩阵可逆，作适当的像坐标置换后，
可以使各阶顺序主子式在该点都非零。
这是带行选主元的消元法的直接结果。
由连续性，缩小邻域后，这些主子式在邻域内都非零。
定义
\[
S_0(x)=x,\qquad
S_k(x)=\bigl(y_1(x),\ldots,y_k(x),x_{k+1},\ldots,x_n\bigr),
\quad 1\le k\le n.
\]
$S_k$ 的 Jacobian 恰为相应的 $k$ 阶顺序主子式。
由反函数定理，可以取同一个充分小的邻域，使所有 $S_k$ 都有连续可导的逆。
令
\[
H_k=S_k\circ S_{k-1}^{-1}.
\]
$H_k$ 只改变第 $k$ 个坐标，而且 Jacobian 非零。
又有
\[
T=H_n\circ\cdots\circ H_1.
\]
依次使用单坐标变换公式，Jacobian 的绝对值因子相乘，
由链式法则正好得到 $|\det DT|$。
这证明了局部公式。

最后用紧集 $\Omega$ 的有限开覆盖，取足够细的正方体网格，
使每个小块都包含于一个分解邻域。
各小块的公共边界及其像均为零体积集，
对各局部公式求和，便得到整体公式。
\end{proof}

\subsubsection*{柱面与球面坐标变换}

在三维空间中，柱面坐标变换为
\[
x=r\cos\theta,\qquad y=r\sin\theta,\qquad z=z,
\]
其中 $0\le r<+\infty$、$0\le\theta\le2\pi$、$-\infty<z<+\infty$。
计算得
\[
\frac{\partial(x,y,z)}{\partial(r,\theta,z)}
=\begin{vmatrix}
\cos\theta&-r\sin\theta&0\\
\sin\theta&r\cos\theta&0\\
0&0&1
\end{vmatrix}=r.
\]
因此相应的体积元素为 $r\,\dd r\,\dd\theta\,\dd z$。

\begin{center}
\begin{tikzpicture}[>=Stealth,font=\small,
  x={(-.8cm,-.55cm)},y={(1cm,0cm)},z={(0cm,1cm)}]
  \draw[->] (0,0,0)--(2.2,0,0) node[below] {$x$};
  \draw[->] (0,0,0)--(0,2.7,0) node[right] {$y$};
  \draw[->] (0,0,0)--(0,0,2.4) node[above] {$z$};
  \node[above left] at (0,0,0) {$O$};
  \draw (0,0,0)--(1.15,1.8,0)--(1.15,1.8,1.7)--cycle;
  \draw[dashed] (1.15,1.8,0)--(1.15,0,0);
  \node[right] at (1.15,1.8,1.7) {$(x,y,z)$};
  \node[right] at (1.15,1.8,0) {$(x,y,0)$};
  \node[below] at (.55,.95,0) {$r$};
  \draw[domain=0:57,samples=24,variable=\t]
    plot ({.6*cos(\t)},{.6*sin(\t)},0);
  \node[below] at (.62,.26,0) {$\theta$};
  \begin{scope}[xshift=5.7cm]
    \draw[->] (0,0,0)--(2.2,0,0) node[below] {$x$};
    \draw[->] (0,0,0)--(0,2.7,0) node[right] {$y$};
    \draw[->] (0,0,0)--(0,0,2.4) node[above] {$z$};
    \node[above left] at (0,0,0) {$O$};
    \draw (0,0,0)--(1.15,1.8,0)--(1.15,1.8,1.7)--cycle;
    \draw[dashed] (1.15,1.8,0)--(1.15,0,0);
    \node[right] at (1.15,1.8,1.7) {$(x,y,z)$};
    \node[right] at (1.15,1.8,0) {$(x,y,0)$};
    \node[above left] at (.6,.94,.9) {$r$};
    \draw[domain=0:57,samples=24,variable=\t]
      plot ({.6*cos(\t)},{.6*sin(\t)},0);
    \node[below] at (.62,.26,0) {$\theta$};
    \draw[domain=0:51.5,samples=24,variable=\t]
      plot ({.62*.538*sin(\t)},{.62*.843*sin(\t)},{.62*cos(\t)});
    \node[right] at (.12,.26,.53) {$\varphi$};
  \end{scope}
\end{tikzpicture}

\end{center}



球面坐标变换为
\[
x=r\sin\varphi\cos\theta,\qquad
y=r\sin\varphi\sin\theta,\qquad z=r\cos\varphi,
\]
其中 $0\le r<+\infty$、$0\le\varphi\le\pi$、$0\le\theta\le2\pi$。
注意 Jacobian 的分母按 $(r,\varphi,\theta)$ 排列：
\[
\begin{aligned}
\frac{\partial(x,y,z)}{\partial(r,\varphi,\theta)}
&=\begin{vmatrix}
\sin\varphi\cos\theta&r\cos\varphi\cos\theta&-r\sin\varphi\sin\theta\\
\sin\varphi\sin\theta&r\cos\varphi\sin\theta&r\sin\varphi\cos\theta\\
\cos\varphi&-r\sin\varphi&0
\end{vmatrix}\\
&=r^2\sin\varphi.
\end{aligned}
\]
所以体积元素为 $r^2\sin\varphi\,\dd r\,\dd\varphi\,\dd\theta$。
若将分母次序写成 $(r,\theta,\varphi)$，行列式变号，但其绝对值不变。

与平面的极坐标一样，上面的完整参数范围包含坐标退化点或角度接缝。
严格应用定理时，应先限制在 $r>0$、$0<\theta<2\pi$，
球面坐标还需 $0<\varphi<\pi$ 的开参数域内。
对有界区域，可以先避开接缝及退化点应用公式，再用零体积集的邻域估计取极限。

\subsubsection*{$n$ 维球面坐标变换}

在 $\R^n$（$n\ge3$）中，球面坐标可写为
\[
\begin{aligned}
x_1&=r\cos\varphi_1,\\
x_k&=r\sin\varphi_1\cdots\sin\varphi_{k-1}\cos\varphi_k,
\qquad 2\le k\le n-1,\\
x_n&=r\sin\varphi_1\cdots\sin\varphi_{n-2}\sin\varphi_{n-1},
\end{aligned}
\]
其中
\[
0\le r<+\infty,\qquad
0\le\varphi_1,\ldots,\varphi_{n-2}\le\pi,\qquad
0\le\varphi_{n-1}\le2\pi.
\]
逐次使用 $\cos^2\alpha+\sin^2\alpha=1$，得到
\[
x_1^2+\cdots+x_n^2=r^2.
\]
相应 Jacobian 的绝对值为
\[
\left|\frac{\partial(x_1,\ldots,x_n)}
{\partial(r,\varphi_1,\ldots,\varphi_{n-1})}\right|
=r^{n-1}\sin^{n-2}\varphi_1\sin^{n-3}\varphi_2\cdots\sin\varphi_{n-2}.
\]

% !TEX root = ../../main.tex
% 来源：Page1.png 左栏、右栏，再接 Page2.png 左栏；右栏为空。
